\hypertarget{cont6:paper-ix-weighted-nonlinear-lifting}{%
\subsection{Chapter IX: weighted nonlinear
lifting}\label{cont6:paper-ix-weighted-nonlinear-lifting}}

\hypertarget{cont6:the-theorem}{%
\subsubsection{The theorem}\label{cont6:the-theorem}}

Let \(p\) be prime, \(R=\mathbb Z_p\), \(K=\mathbb Q_p\), and
\(d,s,r\in\mathbb N\). Let
\(F=(F_1,\ldots,F_d)\in R[X_1,\ldots,X_d]^d\), \(a\in R^d\), and let \[
J=DF(a)=\bigl(\partial_jF_i(a)\bigr)_{i,j}.
\] Assume \[
J\in\operatorname{GL}_d(K),\qquad p^sJ^{-1}\in M_d(R),\qquad
F(a)\in p^rR^d,\qquad r>2s.
\] Then \[
\boxed{\exists!\,b\in a+p^{s+1}R^d:\ F(b)=0,}
\qquad
\boxed{b-a\in p^{r-s}R^d.}
\] All vector divisibility assertions are coordinatewise. The uniqueness
predicate is the entire ball of depth \(s+1\): competitors need not be
assumed to have depth \(r-s\). The theorem permits arbitrary integral
p-adic coefficients and centers, every prime including two, and
dimension zero. It retains the supplied inverse-denominator exponent
\(s\); there is no replacement by a determinant valuation. Integrality
of \(p^sJ^{-1}\) is expressed in Lean by existence of an integral matrix
with exactly that image under the coefficient inclusion
\(R\hookrightarrow K\).

\hypertarget{cont6:proof}{%
\subsubsection{Proof}\label{cont6:proof}}

Write \(A=p^sJ^{-1}\in M_d(R)\). Then \[
AJ=JA=p^sI.
\] The formal matrix bridge proves the equivalence with the stated
fraction-field hypotheses, including nonsingularity. In particular,
multiplication by \(A\) is injective, since multiplication by \(J\)
sends \(Av=0\) to \(p^sv=0\).

Put \(G=AF\), \(t=p^{s+1}\), and \(k=r-(2s+1)\). Taylor expansion of the
actual preconditioned polynomials gives one polynomial system \(Q\) such
that \[
G(a+tu)=G(a)+p^{2s+1}u+p^{2s+2}Q(u),\qquad Q(0)=0.
\] This expansion is proved first over every commutative ring. Its
remainder is a polynomial valid for every \(u\), not a pointwise
quotient depending on \(u\). Since \(G(a)\in p^rR^d\), it can be written
\(p^{2s+1}C\) with \(C\in p^kR^d\). Thus \[
G(a+tu)=p^{2s+1}\bigl(C+u+pQ(u)\bigr).
\]

Integral polynomial evaluation is 1-Lipschitz in the supremum p-adic
metric. Consequently \(T(u)=-C-pQ(u)\) has contraction constant
\(|p|_p<1\) on the complete space \(R^d\), and has exactly one fixed
point \(u\).

For completeness, the sharp precision is additional to this fixed-point
existence. Inductively, if \(u\in p^nR^d\), then polynomial congruence
and \(Q(0)=0\) imply \(Q(u)\in p^nR^d\). For \(n<k\), the equation
\(u=-C-pQ(u)\) therefore gives \(u\in p^{n+1}R^d\). Starting at \(n=0\),
this proves \(u\in p^kR^d\). Hence \(b=a+tu\) satisfies \[
b-a\in p^{s+1+k}R^d=p^{r-s}R^d.
\] The normalization and injectivity of \(A\) give \(F(b)=0\).

Every other \(z\in a+p^{s+1}R^d\) has the form \(a+tw\). If \(F(z)=0\),
the same normalization, with cancellation of the nonzero scalar
\(p^{2s+1}\), gives \(w=T(w)\). Fixed-point uniqueness implies \(w=u\),
hence \(z=b\).

\hypertarget{cont6:boundary-and-sharpness}{%
\subsubsection{Boundary and sharpness}\label{cont6:boundary-and-sharpness}}

The strict inequality cannot in general be weakened to \(r\ge2s\): for
\(p=2\), \(F(X)=X^2+3\), and \(a=1\), we have \(F(a)=4\), \(DF(a)=2\),
and may take \(s=1,r=2\). Yet no root exists in \(\mathbb Z_2\), since
there is none modulo eight. The auxiliary Lean controls prove this
obstruction and the exact residual depth.

The displacement exponent cannot uniformly be increased by one. For
\(F(X)=p^sX+p^r\) at \(a=0\), with \(r>2s\), its root is \(b=-p^{r-s}\),
which is not divisible by \(p^{r-s+1}\). The controls prove the linear
root and the absence of that extra digit for all primes.
